\section{Metaphysics of Dueling}

In De Monarchia, \textbf{Dante}'s goal is to demonstrate that, in general, the power of the Emperor derives from God, and, more specifically, the Roman Empire was such by right. It is instructive to follow Dante's reasoning, as it challenges both the modern way of thinking, such as it is, as well as common misconceptions among those who identify themselves as Traditionalists. If he is a contender as the ``greatest of Europeans'', as Coomaraswamy claims, then anyone whose avowed goal is to ``save Western Civilization'' is duty and honour bound to understand Dante, to know what he is trying to save.

For Dante, the \textbf{Logos} created the world, which is therefore logical, reasonable, and orderly. Man is a \textbf{Rational Soul} who has the power to understand the Logos and hence its Will. Dante takes pains to demonstrate all his points with a logical argument. To buttress his points, Dante refers to secular and religious, Christian and Pagan, sources, as suit his purpose.

The fundamental first point to ponder is whether God, or the gods, are a force in the world. As \textbf{Rene Guenon} insists, everything that happens, does so necessarily and for a reason. The modern mind doubts that, believing instead that things occur randomly with no guiding purpose, in the mistaken belief that this is the more hard-nosed and logical way of thinking. To the contrary, to know something is to know its sufficient reason; that is the very definition of rationality.

The Ancients referred to the will of the gods as Fortune, ``that agency which we more wisely and rightly name Divine Providence,'' as Dante claims. This aligns him more with Guenon, who in The Great Triad wrote of the interactions of Providence, Will, and Destiny\footnote{\url{https://www.gornahoor.net/?p=1798}}. Dante, therefore, takes pains to explain how to discern the Will of God, or Providence. He writes:

\begin{quotex}
For hunting down adequately the truth of our inquiry, it is essential to know that Divine judgment in human affairs is sometimes manifest to men, and sometimes hidden. And it may be manifested in two ways, namely, by reason and by faith.
\end{quotex}


Keep in mind that Reason is used in a wider sense than what is common today, as it includes the Practical Reason: duty, ethics, virtue, etc. Hence, as an example of knowing God's will by Reason, Dante provides examples from \textbf{Aristotle}. The Natural Law is itself Providential, as God is the Author of Nature. Some truths, for Dante, cannot be known by Reason alone, and for these he refers to Scripture. Nevertheless, there are some judgments of God that are hidden and inaccessible both to Reason and Revelation. These cases require a Special Grace to be known. These appear in two ways (along with various ways of manifesting):

\begin{enumerate}
\item \textbf{Simple revelation}
  \begin{itemize}
	\item Spontaneous Act of God
    \begin{itemize}
		\item Expressed directly
		\item Expressed by a sign
    \end{itemize}
	\item Answer to Prayer
  \end{itemize}
\item \textbf{Judicial award}
  \begin{itemize}
	\item By lot
	\item By contest
	\begin{itemize}
		\item Duel, or contest of champions
		\item The contest of athletes
    \end{itemize}
\end{itemize}
\end{enumerate}


It is not my intent at this time to dwell on the simple revelation, so I have simply summarized the ways that is may manifest itself. However, we should point out here that Dante's view in this regard is no different from the Ancients. In our series on the Ancient City, we documented how important actions were decided by prayer and signs or omens. \textbf{Constantine}'s victory was attended by a sign; to disbelieve this is a rather arbitrary decision.

Similarly, to make a decision by lot has an ancient heritage. Dante points out that the thirteenth Apostle to replace Judas was chosen by lot. (To be clear, it was not totally random, since the choice was between equally qualified candidates.) For Christians today to deny the discernment of God's will by drawing lots is also an arbitrary decision. It is well worth pondering why many are discomfited by this. Perhaps because it seems unjust, but Dante makes clear that this is the last resort when the just result cannot be determined by Reason or Revelation. Yet if everything happens for a reason, that is, not by accident, then to what ought we attribute the winning of a lot? To be clear, this has nothing to do with gambling per se, where the players are interested in personal gain, only where there is a disinterested search for justice.

Dante points out that the word contest is derived from certare, ``to make certain''. As an example of the duel, Dante points to the story of \textbf{Hercules} and \textbf{Antaeus}, as told by \textbf{Lucan} and \textbf{Ovid}. As for the second, the example is the story of \textbf{Atalanta} and \textbf{Hippomenes}, again from Ovid. In the latter story Atalanta determines her husband from a footrace. That the winner is revealed by God's will is evident from the aid that Aphrodite provided to Hippomenes. This is how Dante justifies the duel:

\begin{quotex}
Whatever is acquired by single combat duel is acquired with Right. For when human judgment fails, either because it is wrapped in the darkness of ignorance or because it has not the aid of a judge, then, lest judgment should remain forsaken, recourse must be had to Him who so loved her that, by the shedding of His own blood, He met her full demands in death. Hence the Psalm: ``The righteous Lord loveth righteousness.'' This end is accomplished when, with the free consent of the participants, in love and not in hatred of justice, the judgment of God is sought through a mutual trial of bodily and spiritual strength. Because it was first used in single combat of man to man, this trial of strength we call the duel.
\end{quotex}


However, the duel must be the last resort:

\begin{quotex}
But always in quarrels threatening to become matters of war, every effort should be made to settle the dispute through conference, and only as a last resort through battle. \textbf{Tully} and \textbf{Vegetius} both advance this opinion, the former in Moral Duties, and the latter in his book on The Art of War. And as in medical treatment everything is tried before final recourse is had to the knife or fire, so when we have exhausted all other ways of obtaining judgment in a dispute, we may finally turn to this remedy by single combat, compelled thereto by the necessity of justice.
\end{quotex}


That is the first rule of combat. The second is that the contestants enter the combat not in hatred, but in the ``pure zeal for justice by common consent.'' Dante claims that under these conditions, they duelists meet in the name of God, hence God is present with them (according to the Gospels). And if God is present, then it would be a sin to deny that Justice can fail. If that is not a convincing argument, Dante points out that the Ancients believed the same thing, for which he uses the story of \textbf{Pyrrhus}' message to the Romans:

\begin{quotex}
I demand no gold, nor shall you render me a price; we are not barterers in war, but fighters; with steel, not with gold, let each decide the issue of life. Whether \textbf{Hera} wills that you or I shall reign, or whatever fate may bring, let us determine by prowess.
\end{quotex}


On the other hand, a duel should not be settled for a price which would be an injustice as God would not be present. To the objection that one party may be stronger than the other, Dante refers us to David and Goliath. Hence the conclusion: \emph{whatever is acquired by a duel, is acquired with Right and Justice}.

\textit{Homework. Try to discern the workings of Providence in daily life. Act justly. Trial by lot or combat are now forbidden. Why? How much faith or trust in the Will of God, or the gods, would be required to stake everything on the outcome of the drawing of lots or a duel? Men were willing to die trying to be the husband of Atalanta.}


\flrightit{Posted on 2012-01-10 by Cologero}

\begin{center}* * *\end{center}

\begin{footnotesize}
\begin{sffamily}
\texttt{logres on 2012-01-10 at 02:25 said:}

Where two or three are gathered in my name... 

\hfill

\texttt{Boreas on 2012-01-10 at 06:13 said:}

``Homework. Try to discern the workings of Providence in daily life. Act justly.''

synchro: I just forbade my dear daughter an ice cream dessert since she did not follow the rules of dining (sitting in the chair etc.) even after three warnings. I consider this as a just act. I gave my wife a possibility for mercy upon the matter, but as a good parent she declined although she'd wanted the dessert by herself also, yet couldn't have it since it would have caused trouble and arguments with a little girl in a non-rational emotional swing.

``Trial by lot or combat are now forbidden. Why?''

Since the Word of Logos is a penetrating sword and Justice uses the scales instead of a lot-cup.

\hfill

\texttt{HOO on 2012-01-10 at 11:01 said:}

``Heroin Addict Prays to God for a Miracle -- Watch What Happens''\\
\textit{http://www.youtube.com/watch?v=Vn7\_zJEJqww}\footnote{\url{https://www.youtube.com/watch?v=Vn7\_zJEJqww}}


\hfill

\texttt{Cologero on 2012-01-10 at 12:52 said:}

As Dante pointed out, Justice is often hidden, hence the use of the lot or combat. Do we really believe that all questions of justice have been settled and we no longer need to rely on Fortune or combat? That is hardly obvious, which brings us back to Dante's contention.

Hoo, no one is praying for a miracle, rather for justice. Isn't addiction the just result of taking heroin?

\hfill

\texttt{HOO on 2012-01-10 at 14:39 said:}

miracle:

1. An event that appears inexplicable by the laws of nature and so is held to be supernatural in origin or an act of God.

2. One that excites admiring awe.

mid-12c., from O.Fr. miracle, from L. miraculum ``object of wonder'' (in Church L., ``marvelous event caused by God''), from mirari ``to wonder at,'' from mirus ``wonderful,'' from *smeiros, from PIE *(s)mei- ``to smile, be astonished'' (cf. Skt. smerah ``smiling,'' Gk. meidan ``to smile,'' O.C.S. smejo ``to laugh;'' see smile). Replaced O.E. wundortacen, wundorweorc. The Greek words rendered as miracle in the English bibles were semeion ``sign,'' teras ``wonder,'' and dynamis ``power,'' in Vulgate translated respectively as signum, prodigium, and virtus.

You were writing of Spontaneous Acts of God. I wait for a miracle.

And no heroin isn't addictive. However, frequent use can develop into severe psychological dependence and then addiction. As Jung said in this direction, all non-spiritual highs are a search for a replacement for spiritual highs (experiences), as indeed the search for spiritual ``truth'' is a replacement for knowledge of the veracit rites (actions that actualize experinces).

\hfill

\texttt{Boreas on 2012-01-10 at 15:05 said:}

O' Fortuna, velut luna!

Fates (Urd, Verdandi \& Skuld) and Fortuna decide and influence apparent co-incidences according to the hidden will of God. Justice, on the other hand, is many times blind and thus hidden. Combat and dueling is for men of honor in the case that other means don't suffice, usually influenced by Fates and the hidden will, and manifesting the call for Justice. ``Luck'' in the case of Victory is a hidden lustration and grace, apparently called out by Fate \& a minifestation of supreme Justice.

The just result for taking heroin is intoxication.

Last but not the least, sorry about the long link, but this is a must:

\textit{Whether God plays dice}\footnote{\url{https://books.google.com/books?id=JRWgUwJOj6oC&pg=PA95&dq=wolfgang+smith+god+plays+dice&hl=en&sa=X&ei=tgMNT6jIAoittgeOi-yvBQ&ved=0CDUQ6AEwAA\#v=onepage&q&f=false}}

\hfill

\texttt{Cologero on 2012-01-10 at 22:58 said:}

Missing google page from whether God plays dice: summary of main points.

\begin{itemize}


\item quantities pertain to matter, qualities are indicative of essence or form

\item space corresponds to the material and time to the formal aspect of the space-time continuum

\item Prof Smith mentions the male-female complementarity as another example of hylomorphic polarities. We can point out that Guenon likens this to the Purusha-Prakriti teachings of the Vedanta schools.

\item The next topic is yin-yang. Yin corresponds to matter, yang to form.

\item yang is black and represents the unintelligible aspect of the thing or phenomenon. Yang refers to the intelligible aspect.

\item but why, then, the black circle in the white field and vice versa? More is involved than complementarity. There must be a mutual indwelling (perichoresis). Prof Smith then claims this is the key to the problem of determinism.
\end{itemize}


\hfill

\texttt{Cologero on 2012-01-10 at 23:16 said:}

Prof. Smith brings up good points.

So, Boreas, in the midst of our best laid plans, there arises Fortune or Luck, the dot of yin in the field of yang. In some circles there is too much emphasis on Solarity, to the detriment of full understanding. Alchemically, the marriage of the Solar and Lunar cannot be avoided.

Prof. Smith's discussion on natura naturata and natura naturans relates back to that strange dream I posted here in regards to the question of whether each moment of time is new or whether it is restrained by the preceding moment.

Jacob Boehme and Nicholas Berdyaev are good commentators on the void or abyss from which creativity flows.

\hfill

\texttt{Cologero on 2012-01-11 at 10:04 said:}

For a contemporary example of the use of lots to make a judicial decision, please read this:

In 2006, American federal judge Gregory Presnell from the Middle District of Florida ordered opposing sides in a lengthy court case to settle a trivial (but lengthily debated) point over the appropriate place for a deposition using the game of rock-paper-scissors. 11 The ruling in Avista Management v. Wausau Underwriters stated:

Upon consideration of the Motion -- the latest in a series of Gordian knots that the parties have been unable to untangle without enlisting the assistance of the federal courts -- it is ORDERED that said Motion is DENIED. Instead, the Court will fashion a new form of alternative dispute resolution, to wit: at 4:00 P.M. on Friday, June 30, 2006, counsel shall convene at a neutral site agreeable to both parties. If counsel cannot agree on a neutral site, they shall meet on the front steps of the Sam M. Gibbons U.S. Courthouse, 801 North Florida Ave., Tampa, Florida 33602. Each lawyer shall be entitled to be accompanied by one paralegal who shall act as an attendant and witness. At that time and location, counsel shall engage in one game of ``rock, paper, scissors.'' The winner of this engagement shall be entitled to select the location for the 30(b)(6) deposition to be held somewhere in Hillsborough County during the period July 11--12, 2006. 12 

source:
\url{http://en.wikipedia.org/wiki/Rock-paper-scissors}


\hfill

\texttt{Boreas on 2012-01-12 at 12:35 said:}

Thanks for the summary and modification of the link, Cologero.

I must say that the miraculous healing-conversion of the heroin addict and the judicial decision in the comments of this post are truly American.

\hfill

\texttt{Charlotte Cowell on 2012-01-12 at 13:10 said:}

I thought his humility was very touching -- such a small thing but the degree of gratitude shows how desperate he was feelig. Maybe not a full on miracle but certainly a sign and encouragement, such little things are what can get a person through a very hard time... .

\hfill

\texttt{Boreas on 2012-01-12 at 18:02 said:}

Yea, I meant no critique for the man's emotional response or following actions either.

\hfill

\texttt{Charlotte Cowell on 2012-01-12 at 18:31 said:}

no i didn't think you did Boreas, I just wanted to comment on the video as I watched it the other night but didn't say anything at the time


\hfill

\end{sffamily}
\end{footnotesize}

